resumen
\palabrasclave{Predicción de demanda, gestión de inventarios, sistemas de apoyo a la decisión, series de tiempo}

El presente proyecto tiene como objetivo desarrollar un sistema de apoyo a la decisión orientado a la predicción de la demanda de productos en los inventarios.

Esta problemática surge a partir de la falta de información precisa sobre las cantidades necesarias de productos, lo cual dificulta la correcta planificación y puede llevar a que halla desabastecimiento, afectando negativamente la atención de los clientes.

Para abordar esta situación, en este trabajo se propone el uso de modelos de predicción basados en datos históricos de consumo. El sistema permitirá estimar las cantidades requeridas de productos cada mes, así como generar indicadores que apoyen la toma de decisiones para los encargados del manejo del invertario, tales como puntos de reposición y alertas. 

Metodológicamente, se emplean diferentes tecnicas y enfoques de la ingenieria de sistemas para obtener esta  herramienta funcional, que contribuya a mejorar la gestión del inventario, optimizar la planificación de compras y reducir el riesgo de desabastecimiento.